%
% exemplo genérico de uso da classe iiufrgs.cls
% $Id: iiufrgs.tex,v 1.1.1.1 2005/01/18 23:54:42 avila Exp $
%
% This is an example file and is hereby explicitly put in the
% public domain.
%
% Ensure url is loaded with hyphenation support before hyperref
\PassOptionsToPackage{hyphens}{url}
\documentclass[ppgc,prop-tese,english]{iiufrgs}
% Para usar o modelo, deve-se informar o programa e o tipo de documento.
% Programas :
%   * cic       -- Graduação em Ciência da Computação
%   * ecp       -- Graduação em Ciência da Computação
%   * ppgc      -- Programa de Pós Graduação em Computação
%   * pgmigro   -- Programa de Pós Graduação em Microeletrônica
%
% Tipos de Documento:
%   * tc                -- Trabalhos de Conclusão (apenas cic e ecp)
%   * diss ou mestrado  -- Dissertações de Mestrado (ppgc e pgmicro)
%   * tese ou doutorado -- Teses de Doutorado (ppgc e pgmicro)
%   * ti                -- Trabalho Individual (ppgc e pgmicro)
%
% Outras Opções:
%   * english    -- para textos em inglês
%   * openright  -- Força início de capítulos em páginas ímpares (padrão da
%                   biblioteca)
%   * oneside    -- Desliga frente-e-verso
%   * nominatalocal -- Lê os dados da nominata do arquivo nominatalocal.def


% Use unicode
\usepackage[utf8]{inputenc}   % pacote para acentuação

%using subfigures and tables
\usepackage{caption}
\usepackage{subcaption}

% Necessário para incluir figuras
\usepackage{graphicx}           % pacote para importar figuras


\usepackage{times}              % pacote para usar fonte Adobe Times
% \usepackage{palatino}
% \usepackage{mathptmx}          % p/ usar fonte Adobe Times nas fórmulas

\usepackage[alf,abnt-emphasize=bf]{abntex2cite}	% pacote para usar citações abnt

\usepackage{longtable}
\usepackage{subfiles} % Best loaded last in the preamble
\usepackage[inline]{enumitem}

% For theorems and such
\usepackage{mathtools}
\usepackage{amssymb}
\usepackage{xcolor}
\usepackage{cite}
\usepackage{mathrsfs}
\usepackage{algorithmic}
\usepackage{textcomp}

% \newcommand{\repourl}{https://anonymous.4open.science/r/grokking_explained-2822}
\newcommand{\repourl}{https://github.com/brenowca/grokking_explained}

\newtheorem{definition}{Definition}
\def\BibTeX{{\rm B\kern-.05em{\sc i\kern-.025em b}\kern-.08em
    T\kern-.1667em\lower.7ex\hbox{E}\kern-.125emX}}
    
\DeclarePairedDelimiter{\ceil}{\lceil}{\rceil}
\DeclarePairedDelimiter{\floor}{\lfloor}{\rfloor}
% \newcommand{\centroid}[1]{\mathbf{C}({#1})}
\newcommand{\centroid}[1]{\mathscr{C}({#1})}
% \newcommand{\centroid}[1]{\overline{#1}}


%
% Informações gerais
%
\title{Towards a Local Interpretation of Attention Based Neural Models}
\translatedtitle{Rumo a uma Interpretação Local de Modelos Neurais Baseados em Atençãoz}

\author{Santos Rezende de Carvalho}{Breno William}
% alguns documentos podem ter varios autores:
%\author{Flaumann}{Frida Gutenberg}
%\author{Flaumann}{Klaus Gutenberg}

% orientador e co-orientador são opcionais (não diga isso pra eles :))
\advisor[Prof.~Dr.]{Lamb}{Luís}
\coadvisor[Prof.~Dr.]{d'Avila Garcez}{Artur}

% a data deve ser a da defesa; se nao especificada, são gerados
% mes e ano correntes
%\date{May}{2025}

% o local de realização do trabalho pode ser especificado (ex. para TCs)
% com o comando \location:
%\location{Itaquaquecetuba}{SP}

% itens individuais da nominata podem ser redefinidos com os comandos
% abaixo:
% \renewcommand{\nominataReit}{Prof\textsuperscript{a}.~Wrana Maria Panizzi}
% \renewcommand{\nominataReitname}{Reitora}
% \renewcommand{\nominataPRE}{Prof.~Jos{\'e} Carlos Ferraz Hennemann}
% \renewcommand{\nominataPREname}{Pr{\'o}-Reitor de Ensino}
% \renewcommand{\nominataPRAPG}{Prof\textsuperscript{a}.~Joc{\'e}lia Grazia}
% \renewcommand{\nominataPRAPGname}{Pr{\'o}-Reitora Adjunta de P{\'o}s-Gradua{\c{c}}{\~a}o}
% \renewcommand{\nominataDir}{Prof.~Philippe Olivier Alexandre Navaux}
% \renewcommand{\nominataDirname}{Diretor do Instituto de Inform{\'a}tica}
% \renewcommand{\nominataCoord}{Prof.~Carlos Alberto Heuser}
% \renewcommand{\nominataCoordname}{Coordenador do PPGC}
% \renewcommand{\nominataBibchefe}{Beatriz Regina Bastos Haro}
% \renewcommand{\nominataBibchefename}{Bibliotec{\'a}ria-chefe do Instituto de Inform{\'a}tica}
% \renewcommand{\nominataChefeINA}{Prof.~Jos{\'e} Valdeni de Lima}
% \renewcommand{\nominataChefeINAname}{Chefe do \deptINA}
% \renewcommand{\nominataChefeINT}{Prof.~Leila Ribeiro}
% \renewcommand{\nominataChefeINTname}{Chefe do \deptINT}

%
% Palavras-chave
% A primeira palavra de cada palavra-chave (que, ironicamente, pode ser várias palavras) deve sempre
% começar em maíuscula (novo padrão da biblioteca), as demais palavras começam em minúscula, exceto
% se são nomes próprios.
%
\keyword{Deep Learning interpretability} 
\keyword{Neural-networks interpretability}
\keyword{Artificial Intelligence Alignment}
\keyword{High-dimensional time-series visualization}

\keyword{Attention-based models}
\keyword{Grokking}
\keyword{Sparse autoencoders}
\keyword{Cross-layer transcoders}
\keyword{Neuro-symbolic AI}
\keyword{Logic Tensor Networks}

%
% palavras-chave na lingua estrangeira
% A primeira palavra de cada palavra-chave (que, ironicamente, pode ser várias palavras) deve sempre
% começar em maíuscula (novo padrão da biblioteca), as demais palavras começam em minúscula, exceto
% se são nomes próprios. Nomes próprios não dvem ser traduzidos.
%
\translatedkeyword{Interpretabilidade de Deep Learning}
\translatedkeyword{Interpretabilidade de redes neurais}
\translatedkeyword{Alinhamento de Inteligência Artificial}
\translatedkeyword{Visualização de series temporais de alta dimensionalidade}

\translatedkeyword{Modelos baseados em atenção}
\translatedkeyword{Grokking}
\translatedkeyword{Auto-encoder esparso}
\translatedkeyword{Cross-layer transcoders}
\translatedkeyword{IA neuro-symbólica}
\translatedkeyword{Redes Lógicas de Tensor}
%
% inicio do documento
%
\begin{document}

% folha de rosto
% às vezes é necessário redefinir algum comando logo antes de produzir
% a folha de rosto:
% \renewcommand{\coordname}{Coordenadora do Curso}
\maketitle

% dedicatoria
\clearpage
\subfile{chapters/abstract}

% lista de figuras
\listoffigures

% lista de tabelas
\listoftables

% lista de abreviaturas e siglas
% o parametro deve ser a abreviatura mais longa
% A NBR 14724:2011 estipula que a ordem das abreviações
% na lista deve ser alfabética.
\begin{listofabbrv}{SPMD}
        \item[DL]   Deep Learning
        
        \item[CLT]  Cross-layer Transcoder
        \item[LLM]  Large Language Model
        \item[LTN]  Logic Tensor Networks
        \item[SAE]  Sparse Auto-Encoder
        \item[PCA]  Principal Component Analysis
        
        \item[t-SNE]  t-Distributed Stochastic Neighbor Embedding
\end{listofabbrv}

% idem para a lista de símbolos
%\begin{listofsymbols}{$\alpha\beta\pi\omega$}
%       \item[$\sum{\frac{a}{b}}$] Somatório do produtório
%       \item[$\alpha\beta\pi\omega$] Fator de inconstância do resultado
%\end{listofsymbols}

% sumario
\tableofcontents

% aqui comeca o texto propriamente dito
% \chapter{Instruções}
%     \subfile{chapters/instructions}

%\part{Understanding}
\chapter{Introduction}
    \label{chpt:intro}
    \subfile{chapters/intro}
    
% \chapter{Alignment}
%     \label{chpt:alignment}
%     \subfile{chapters/alignment}

\chapter{Background}
    \label{chpt:background}
    \subfile{chapters/background}

\chapter{Dynamics of Emergence: Training phenomena under distribution shift}%: Using interpretability and statistical tools}
% \chapter{Interesting learning phenomena}
    \label{chpt:learning_phenomena}
    \subfile{chapters/learning_phenomena}

%\part{Steering}
\chapter{Methodology}
    \label{chpt:circuitry}
    \subfile{chapters/circuitry}

\chapter{Current Results}
    \label{chpt:constrained_generation}
    \subfile{chapters/prop-tese/current-results}

%\part{Generality Across Modalities}

\chapter{Timetable}
    \label{chpt:reasoning}
    \subfile{chapters/prop-tese/timetable}

\chapter{Contributions and risks}
    \label{chpt:future}
    \subfile{chapters/prop-tese/future}

\chapter{Conclusion}
    \label{chpt:closing}
    \subfile{chapters/closing}

% e aqui vai a parte principal
%
% \chapter{Estado da arte}
% \chapter{Mais estado da arte}
% \chapter{A minha contribuição}
% \chapter{Prova de que a minha contribuição é válida}
% \chapter{Conclusão}

% referencias
% aqui será usado o environment padrao `thebibliography'; porém, sugere-se
% seriamente o uso de BibTeX e do estilo abnt.bst (veja na página do
% UTUG)
%
% observe também o estilo meio estranho de alguns labels; isso é
% devido ao uso do pacote `natbib', que permite fazer citações de
% autores, ano, e diversas combinações desses

\bibliographystyle{abntex2-alf}
\bibliography{biblio,biblio-grokking}


\appendix

\chapter{Resumo expandido}

\noindent
\textbf{Resolução 02/2021 -- Redação de Teses e Dissertações em Inglês}
Dissertações de Mestrado e Teses de Doutorado do PPGC, bem como outros
trabalhos escritos tais como Proposta de Tese e PEP, poderão ser
redigidas em inglês desde que contenham um título e resumo expandido
redigidos em português. O resumo expandido deve conter no mínimo duas
páginas inteiras, deve aparecer como apêndice e deve conter as
principais contribuições e resultados do trabalho.

\vspace{0.8em}
\noindent
Esta proposta de tese investiga interpretabilidade local em modelos neurais baseados em atenção, com foco em três eixos complementares: (i) dinâmica de aprendizagem e emergência de generalização tardia (grokking), (ii) controle comportamental por restrições neuro-simbólicas, e (iii) diagnóstico de mudanças de distribuição e perturbações adversariais a partir de sinais internos do modelo. O problema central é que modelos modernos apresentam alto desempenho prático, mas ainda carecemos de métodos confiáveis para explicar, em nível local, por que uma saída específica foi produzida, como estruturas internas emergem durante o treinamento e como intervir com baixo efeito colateral.

\section*{Motivação e problema}
Em aplicações reais, desempenho médio não é suficiente: é necessário compreender falhas, antecipar degradações e justificar decisões em casos individuais. Métodos de interpretabilidade pós-hoc frequentemente oferecem explicações plausíveis, porém nem sempre fiéis aos mecanismos internos do modelo. Em contrapartida, abordagens mecanísticas tendem a ser mais rigorosas, mas exigem protocolos de avaliação claros para conectar medições internas a efeitos observáveis no comportamento.

Neste contexto, esta tese parte da hipótese de que comportamentos de alto nível emergem de estruturas internas parcialmente identificáveis (características, circuitos e padrões de ativação) que podem ser medidas, alinhadas entre camadas e usadas para diagnóstico e controle local. A contribuição proposta não é uma promessa de interpretabilidade global completa, mas um arcabouço experimentalmente testável para análise e intervenção local, com métricas explícitas e avaliação de trade-offs.

\section*{Objetivos}
Os objetivos científicos são:
\begin{enumerate}
    \item caracterizar mudanças representacionais ao longo do treinamento, especialmente em transições de generalização tardia;
    \item desenvolver e avaliar um pipeline de steering por restrições lógicas suaves, medindo localidade do efeito e regressões colaterais;
    \item detectar mudanças de distribuição e ataques a partir de assinaturas internas em ativações/atenção, complementando sinais baseados apenas na saída.
\end{enumerate}

Esses objetivos são avaliados por métricas comportamentais (acurácia, aderência a restrições, qualidade de geração, ROC-AUC/PR-AUC em detecção) e métricas mecanísticas (esparsidade, erro de reconstrução, estabilidade entre sementes, alinhamento entre camadas e satisfação de regras lógicas).

\section*{Metodologia proposta}
A metodologia integra três componentes principais:
\begin{itemize}
    \item \textbf{SAEs (Sparse Autoencoders):} extração de códigos esparsos por camada para reduzir polissemia e facilitar interpretação de características latentes;
    \item \textbf{Cross-layer transcoders:} mapeamento entre espaços latentes de diferentes camadas para medir consistência semântica ao longo da profundidade;
    \item \textbf{Validadores lógicos (LTN-style):} predicados e regras diferenciáveis usados tanto como critério de avaliação quanto como sinal de controle no processo de geração.
\end{itemize}

O desenho experimental é organizado em três famílias de tarefas: (i) cenários controlados de grokking e distribuição deslocada; (ii) tarefas de geração para steering por restrições; e (iii) cenários de robustez com corrupção/adversarial. Em cada família, são executadas ablações (remoção de cada módulo do pipeline) e comparações com baselines para isolar contribuição causal dos componentes.

\section*{Resultados já obtidos}
Dois resultados prévios sustentam a viabilidade da proposta. Primeiro, o trabalho \cite{carvalho2025inducing_grokking_distribution_shifts} mostrou que mudanças de distribuição podem induzir grokking de forma sistemática em cenários controlados, com evidências de que estrutura relacional dos dados influencia fortemente a generalização tardia. Segundo, o estudo \cite{carvalho2026interpretable_stego_detection_sae} indicou que representações esparsas em modelos de visão (ViT/CLIP) permitem separar entradas limpas e perturbadas com poucas características latentes interpretáveis, alcançando resultados competitivos de detecção.

Esses resultados prévios são incorporados nesta tese como base metodológica: uso de ambientes controlados para validar hipóteses mecanísticas e uso de assinaturas latentes esparsas para diagnóstico local.

\section*{Contribuições esperadas}
As contribuições esperadas da tese são:
\begin{itemize}
    \item um framework unificado de interpretabilidade local para modelos baseados em atenção, conectando aprendizado, inferência e controle;
    \item evidências empíricas sobre como \textit{sparse features} se reorganizam em fases de generalização tardia;
    \item um protocolo de steering com medição explícita de trade-off entre satisfação de restrições e qualidade de saída;
    \item sinais diagnósticos práticos para detecção precoce de mudança de regime e perturbação adversarial.
\end{itemize}

\section*{Riscos e mitigação}
Os principais riscos incluem instabilidade de \textit{latent features} entre sementes, alinhamento entre camadas sem preservação semântica, degradação de qualidade sob restrições fortes e baixa transferência de assinaturas de ataque para cenários mais realistas. Como mitigação, a proposta prevê múltiplas sementes, avaliação estatística com intervalos de confiança, ablações sistemáticas, calibração de pesos de regras e relato explícito de falhas e limites de generalização.

\section*{Síntese}
Em síntese, esta proposta organiza uma agenda experimental para conectar interpretabilidade mecanística e utilidade prática. O foco em critérios locais, mensuráveis e reprodutíveis busca produzir conhecimento acionável para análise, controle e robustez de modelos neurais baseados em atenção, contribuindo para seu uso mais seguro e previsível em aplicações reais.


% \chapter{Final Remarks}


\end{document}


